\subsection{変数係数の一次漸化式}\label{節：変数係数の一次漸化式}
  \sectionref*{節：階比型}では添字ごとの倍率を掛け合わせ,
  \sectionref*{節：階差等比複合型}では同じ更新をする量を引いた.
  ここでは,倍率と付加項がどちらも添字に依存する場合を,この二つの操作で解こう.

  \subsubsection{倍率をそろえて,変化を足す}
  まず
  \BookMathLayout{\[
    a_1=2,\qquad
    a_{n+1}=\frac{n+2}{n}a_n+(n+1)(n+2)
    \quad(n\ge1)
  \]}{\[
    \begin{aligned}
      a_1&=2,\\
      a_{n+1}&=\frac{n+2}{n}a_n+(n+1)(n+2)\\
      &\quad(n\ge1)
    \end{aligned}
  \]}
  を考える.係数の分子・分母に注目すると,
  \[
    \frac{(n+1)(n+2)}{n(n+1)}=\frac{n+2}{n}
  \]
  である.\textbf{隣り合う項の比が元の倍率になる量}として$h_n=n(n+1)$を選べる.
  これは\sectionref*{節：補助数列の設計}で目標から割る量を選んだ方法である.
  $b_n=\frac{a_n}{n(n+1)}$と置いて実際に計算すると,
  \[
    \frac{a_{n+1}}{(n+1)(n+2)}
      =\frac{a_n}{n(n+1)}+1
  \]
  となる.従って$b_1=1$, $b_{n+1}=b_n+1$より$b_n=n$であり,
  \[
    a_n=n^2(n+1)
  \]
  を得る.分母は$n\ge1$で非零である.初項と元の式へ代入しても確かめられる.

  \subsubsection{一般の手順と,割り算の条件}
  $a_{n+1}=p_na_n+q_n$ $(n\ge1)$に対して,
  $h_1=1$, $h_{n+1}=p_nh_n$とする.
  すべての$p_n$が非零なら,$h_n$も非零であり,
  \[
    h_n=\prod_{k=1}^{n-1}p_k
  \]
  である.空積は$1$とする.両辺を$h_{n+1}$で割ると,
  \[
    b_n=\frac{a_n}{h_n},\qquad
    b_{n+1}-b_n=\frac{q_n}{h_{n+1}}
  \]
  だから,階差を足して
  \[
    a_n=h_n\left(a_1+\sum_{k=1}^{n-1}\frac{q_k}{h_{k+1}}\right)
  \]
  を得る.先ほどは定数倍して$h_n=n(n+1)$を選んだが,
  その場合の$b_1$を正しく合わせれば同じ答えになる.

  \begin{bookpoint}
    \textbf{正規化の役割}\par
    倍率を取り除くと,残った付加項を足す問題になる.
    積・和の表示を得ることと,最後まで簡単な式で計算できることは別である.
  \end{bookpoint}
  もし$p_j=0$なら,$a_{j+1}=q_j$はそれ以前の項によらず決まる.
  零を含む$h_n$で割らず,零となる更新の前後で区間を分け,
  非零の倍率が続く範囲で同じ方法を使う.

  \subsubsection{同じ更新をする量を引く別解}
  先ほどの例では,$v_n=n^2(n+1)$が元の更新を満たすことを直接確かめられる.
  この候補は,正規化後の$b_n=n$から見つけたものである.
  $d_n=a_n-v_n$と置くと,
  \[
    d_{n+1}=\frac{n+2}{n}d_n
  \]
  となり,階比型へ戻る.初項を一般の$a_1=s$に変えると,
  \[
    a_n=n(n+1)\left(n+\frac{s-2}{2}\right)
  \]
  である.一つの解$v_n$と,同次の解の定数倍に分かれている.
  これは\sectionref*{節：非同次と解の組合せ}の見方そのものである.

  \begin{bookpoint}
    \textbf{二つの方法の選び方}\par
    隣り合う比が分かるなら割る因子を探す.
    同じ更新をする既知の量が見えるならそれを引く.
    どちらも,付加項と倍率を分けて扱うための操作である.
  \end{bookpoint}
  操作だけの比較は\sectionref*{節：操作を選ぶ練習},
  一般項まで求める練習は\bookproblemref{practice-connections}{1}{接続演習問1}へ進もう.
